\begingroup
\renewcommand{\arraystretch}{1.45}
\Needspace{23\baselineskip}
\InterpolationArticle{13}
\textsc{Satz.}\quad\textit{Ist $X$ eine Funktion des Bogens $x$ von der Form $(F)$}
\[
\begin{aligned}
&\alpha+\alpha'\cos x+\alpha''\cos 2x+\alpha'''\cos 3x+\text{usw.}\\[5pt]
&\phantom{\alpha}+\beta'\sin x+\beta''\sin 2x+\beta'''\sin 3x+\text{usw.}
\end{aligned}
\]
\textit{oder von der Form $(G)$}
\[
\begin{aligned}
&\gamma\cos\tfrac12x+\gamma'\cos\tfrac32x+\gamma''\cos\tfrac52x+\text{usw.}\\[5pt]
&+\delta\sin\tfrac12x+\delta'\sin\tfrac32x+\delta''\sin\tfrac52x+\text{usw.}
\end{aligned}
\]
\textit{und ist nach Einsetzung von $x=a$ der Wert der Funktion $X$ $=0$, so ist $X$ durch $\sin\tfrac12(x-a)$ teilbar, und der Quotient hat im ersten Fall die Form $G$, im zweiten die Form $F$.}

\textit{Beweis.\quad Erster Fall.}\quad Ersetzt man in der Funktion $X$ jedes Glied $\cos nx$ durch $\cos nx-\cos na$, jedes Glied $\sin nx$ hingegen durch $\sin nx-\sin na$,
\SourceLeaf{296}
und schließlich $\alpha$ durch $\alpha-\alpha$ oder $0$, so heben sich die auf diese Weise hinzugefügten konstanten Glieder offenbar gegenseitig auf, sodass die Funktion $X$ unverändert bleibt. Jedes einzelne Glied von $X$ ist nun durch $\sin\tfrac12(x-a)$ teilbar, nämlich
\[
\begin{aligned}
\frac{\cos nx-\cos na}{\sin\tfrac12(x-a)}
&=-2\sin((n-\tfrac12)x+\tfrac12a)\\[5pt]
&\phantom{=}{-}2\sin((n-\tfrac32)x+\tfrac32a)\\[5pt]
&\phantom{=}{-}2\sin((n-\tfrac52)x+\tfrac52a)\\[5pt]
&\phantom{=}{-}\text{usw.}\\[5pt]
&\phantom{=}{-}2\sin((\tfrac12x+(n-\tfrac12)a)
\end{aligned}
\]
\[
\begin{aligned}
\frac{\sin nx-\sin na}{\sin\tfrac12(x-a)}
&=+2\cos((n-\tfrac12)x+\tfrac12a)\\[5pt]
&\phantom{=}{+}2\cos((n-\tfrac32)x+\tfrac32a)\\[5pt]
&\phantom{=}{+}2\cos((n-\tfrac52)x-\tfrac52a)\\[5pt]
&\phantom{=}{+}\text{usw.}\\[5pt]
&\phantom{=}{+}2\cos(\tfrac12x+(n-\tfrac12)a)
\end{aligned}
\]
Beide Koeffizienten lassen sich offenbar auf die Form $G$ bringen; daher ist auch die ganze Funktion $X$ durch $\sin\tfrac12(x-a)$ teilbar, und der Quotient lässt sich auf die Form $G$ bringen. Q.~E.~D.

\Needspace{40\baselineskip}
\textit{Zweiter Fall.}\quad Ersetzt man in der Funktion $X$ jedes Glied $\cos nx$ (wobei $n$ nun nicht wie zuvor eine ganze Zahl, sondern die Hälfte einer ungeraden ganzen Zahl bezeichnet) durch $\cos nx-\cos\tfrac12(x-a)\cos na$, jedes Glied $\sin nx$ hingegen durch $\sin nx-\cos\tfrac12(x-a)\sin na$, so heben sich die Glieder, die auf diese Weise zur Funktion $X$ hinzukommen, offenbar gegenseitig auf, sodass $X$ unverändert bleibt. Jedes einzelne Glied von $X$ ist nun durch $\sin\tfrac12(x-a)$ teilbar, nämlich
\[
\begin{aligned}
\frac{\cos nx-\cos\tfrac12(x-a)\cos na}{\sin\tfrac12(x-a)}\qquad\\[5pt]&=-2\sin((n-\tfrac12)x+\tfrac12a)\\[5pt]
&\phantom{=}{-}2\sin((n-\tfrac32)x+\tfrac32a)\\[5pt]
&\phantom{=}{-}2\sin((n-\tfrac52)x+\tfrac52a)\\[5pt]
&\phantom{=}{-}\text{usw.}\\[5pt]
&\phantom{=}{-}2\sin(x+(n-1)a)\\[5pt]
&\phantom{=}{-}\sin na
\end{aligned}
\]
\[
\begin{aligned}
\frac{\sin nx-\cos\tfrac12(x-a)\sin na}{\sin\tfrac12(x-a)}\qquad\\[5pt]&=+2\cos((n-\tfrac12)x+\tfrac12a)\\[5pt]
&\phantom{=}{+}2\cos((n-\tfrac32)x+\tfrac32a)\\[5pt]
&\phantom{=}{+}2\cos((n-\tfrac52)x+\tfrac52a)\\[5pt]
&\phantom{=}{+}\text{usw.}\\[5pt]
&\phantom{=}{+}2\cos(x+(n-1)a)\\[5pt]
&\phantom{=}{+}\cos na
\end{aligned}
\]
\SourceLeaf{297}
wie sich durch Multiplikation leicht bestätigt. Beide Quotienten haben die Form $F$. Daher ist die ganze Funktion $X$ durch $\sin\tfrac12(x-a)$ teilbar, und der Quotient lässt sich auf die Form $F$ bringen. Q.~E.~D.

\Needspace{32\baselineskip}
\InterpolationArticle{14}
\textsc{Satz.}\quad\textit{Wird eine Funktion $X$ von der Form $F$ oder $G$ für mehrere verschiedene\HistoricalNote{24} Werte von $x$, nämlich $a$, $b$, $c$, $d$ usw., $=0$, so ist sie durch das Produkt}
\[
\sin\tfrac12(x-a)\sin\tfrac12(x-b)\sin\tfrac12(x-c)\sin\tfrac12(x-d)\ .\,.\,.\,.
\]
\textit{teilbar, und der Quotient hat dieselbe Form wie $X$ oder die andere Form, je nachdem, ob die Anzahl der Werte $a$, $b$, $c$, $d$ $.\,.\,.\,.$ gerade oder ungerade ist.}

\textit{Beweis.}\quad Man setze $X=X'\sin\tfrac12(x-a)$, und $M$ sei der Wert von $X'$ für $x=b$. Dann ist $M\sin\tfrac12(b-a)$ der Wert von $X$ für $x=b$; da dieser $=0$ sein muss, gilt notwendig $M=0$. Also ist $X'$ durch $\sin\tfrac12(x-b)$ teilbar und aus demselben Grund der daraus entstehende Quotient durch $\sin\tfrac12(x-c)$, und so weiter. Daher ist $X$ teilbar durch das Produkt
\[
\sin\tfrac12(x-a)\sin\tfrac12(x-b)\sin\tfrac12(x-c)\sin\tfrac12(x-d)\ .\,.\,.\quad\text{Q.~E.~D.}
\]
Der zweite Teil des Satzes ist so offensichtlich, dass kein Beweis nötig ist.

Im Übrigen ist der umgekehrte Satz ebenso offensichtlich: Ist $X$ durch das Produkt
\[
\sin\tfrac12(x-a)\sin\tfrac12(x-b)\sin\tfrac12(x-c)\sin\tfrac12(x-d)\ .\,.\,.\,.
\]
teilbar, so wird sein Wert $=0$, wenn $x$ irgendeiner der Werte $a$, $b$, $c$, $d$ usw. zugewiesen wird.

\endgroup
